\subsubsection{Single layer with imposed temperature b.c.}

Let us take a single layer of material characterised by
a heat capacity $c_p$, a heat conductivity $k$
and a heat production term $H$.

\begin{center}
\includegraphics[width=5cm]{images/initial_temperature/tempcond.png}
\end{center}

The Heat transport equation writes
\[
\rho c_p ( \frac{\partial T}{\partial t} + {\vec v} \cdot {\vec \nabla} { T}) = 
{\vec \nabla} \cdot (k {\vec \nabla} T) + \rho H
\]
At steady state and in the absence of a velocity field, and assuming
that the material properties to be independent of time and space, and that
there is no heat production ($H=0$), this equation
simplifies to
\[
\Delta T =0 
\]
Assuming the layer to be parallel to the $x$-axis, this yields to write
\[
T(x,y)=T(y)=\alpha T+ \beta
\]
In order to specify the constants $\alpha$ and $\beta$, we need two constraints.

At the bottom of the layer $y=y_b$ a temperature $T_b$ is prescribed while a temperature
$T_t$ is prescribed at the top with $y=y_t$. This ultimately yields a temperature field in
the layer given by
\[
\boxed{
T(y) = \frac{T_t-T_b}{y_t-y_b}(y-y_b) + T_b
}
\]

If now the heat production coefficient is not zero, the differential equation
reads
\[
 k \Delta T + H = 0 
\]
The temperature field is then expected to be of the form
\[
T(y)= - \frac{H}{2k} y^2 + \alpha y + \beta 
\]
Supplied again with the same boundary conditions, this leads to
\[
\beta=T_b + \frac{H}{2k} y_b^2 - \alpha y_b
\]
ie,
\[
T(y) = -\frac{H}{2k} (y^2-y_b^2) + \alpha (y-y_b) + T_b
\]
and finally
\[
\alpha =  \frac{T_t-T_b}{y_t-y_b}  + \frac{H}{2k}(y_b+y_t)
\]
or,
\[
T(y) = -\frac{H}{2k} (y^2-y_b^2) + \left( \frac{T_t-T_b}{y_t-y_b}  + \frac{H}{2k}(y_b+y_t)   \right) (y-y_b) + T_b
\]

Taking $H=0$ in this equation obviously yields the temperature field obtained previously.
Taking $k=2.25$, $T_t=0C$, $T_b=550C$, $y_t=660km$, $y_b=630km$ yields the following
temperature profiles and heat fluxes when the heat production $H$ varies:
\begin{center}
\includegraphics[width=5cm]{images/initial_temperature/temperature1.pdf}
\includegraphics[width=5cm]{images/initial_temperature/heatflux1.pdf}
\end{center}
Looking at the values at the top, which are somewhat estimated to be
about $55-65mW/m^2$ \cite[table 8.6]{jama}, one sees that value $H=0.8e-6$ yields a very acceptable
heat flux.
Looking at the bottom, the heat flux is then about $0.03W/m^2$
which is somewhat problematic since the heat flux at the Moho
is reported to be somewhere between 10 and 20 $mW/m^2$ in \cite[table 7.1]{jama}.


%-----------------------------------------------------
\subsection{Single layer with imposed heat flux b.c.}

Let us now assume that heat fluxes are imposed at the top and bottom of the layer:
\begin{center} 
\includegraphics[width=5cm]{images/initial_temperature/tempcond2.png}
\end{center}

We start again from the ODE
\[
k \Delta T + H = 0 
\]
but only integrate it once:
\[
k \frac{dT}{dy}  + H y + \alpha  = 0 
\]
At the bottom $q=k(dT/dy)|_{y=y_b} = q_b$ and at the top
$q=k(dT/dy)|_{y=y_t} = q_t$ so that 

\todo[inline]{to finish}


 
%-----------------------------------------------------
\subsection{Single layer with imposed heat flux and temperature b.c. }

\begin{center}
\includegraphics[width=5cm]{images/initial_temperature/tempcond3.png}
\end{center}

\todo[inline]{to finish}


%---------------------------------------------------------------
\subsubsection{Half cooling space}

%---------------------------------------------------------------
\subsubsection{Plate model}

%---------------------------------------------------------------
\subsubsection{McKenzie slab}










